\section{Problemática}

En las áreas de producción y ensamblaje de TenarisTamsa se utiliza diariamente el Sistema Integral de Producción para el seguimiento y la planificación de los pedidos de las distintas gamas de productos que la empresa provee a sus clientes. Este sistema permite gestionar información relacionada con los procesos productivos, facilitando la organización de pedidos y el monitoreo de las diferentes etapas de fabricación.

No obstante, debido al crecimiento constante de la empresa y a su expansión como proveedor para diversos clientes, los requerimientos operativos dentro del área productiva se han vuelto cada vez más complejos. Entre estos requerimientos se encuentran la necesidad de contar con información actualizada en tiempo real, la trazabilidad completa de los pedidos, la integración de datos entre distintas áreas y plantas, así como la capacidad de analizar información histórica para la toma de decisiones y la mejora continua de los procesos.

Sin embargo, el sistema actualmente utilizado presenta diversas limitaciones que dificultan el cumplimiento de dichos requerimientos operativos. Una de las principales limitaciones es el uso de registros manuales en papel por parte de los trabajadores, lo que incrementa el riesgo de errores, pérdida de información y retrasos en la actualización de los datos. Asimismo, existen múltiples sistemas de información que operan de manera independiente, sin una integración adecuada, lo que obliga a consultar diferentes fuentes por separado y dificulta la consolidación y análisis de la información.

Adicionalmente, al formar parte del grupo corporativo Techint Group, la empresa requiere mantener homogeneidad en los sistemas utilizados por las diferentes plantas del grupo, tales como TenarisSiderca y TenarisDalmine, entre otras, las cuales realizan procesos similares. Este contexto implica como requerimiento operativo la interoperabilidad entre sistemas y la estandarización de la información.

No obstante, los sistemas actuales no permiten una integración eficiente con las plantas previamente mencionadas, lo que limita el intercambio oportuno de información, especialmente en situaciones relacionadas con la mejora de procesos o la prevención de problemas potenciales en otras instalaciones.

En consecuencia, surge la necesidad de analizar y desarrollar soluciones que permitan optimizar la gestión de la información dentro de los procesos productivos, reducir las limitaciones actuales y cumplir con los requerimientos operativos de integración, trazabilidad y disponibilidad de la información, así como garantizar la interoperabilidad de los sistemas entre las distintas plantas del grupo.